\chapter{Sorting and Greedy Algorithms}\label{ch:sorting}

Sorting is the most useful preprocessing step in competitive programming. Once data is ordered, equal values sit next to
each other, extremes sit at the ends, and a whole family of questions collapses into a single left-to-right scan. Greedy
algorithms, which build a solution by repeatedly making the choice that looks best right now, are sorting's natural partner:
almost every greedy algorithm begins by sorting the input into the order in which its choices will be made. This chapter
treats both, with an emphasis on the part that contests actually test: \emph{finding} the right order, and \emph{proving}
that the greedy choice is safe.

\begin{chapterintro}{\bronze\ \silver\ \gold}{Chapters~\ref{ch:approach} to~\ref{ch:ds-basics}, in particular the exchange
argument of Section~\ref{sec:exchange}.}
  \item use sorting to turn questions about pairs and groups into scans, and pick the right library routine;
  \item sort integers in linear time with counting sort and radix sort;
  \item design comparators with exchange arguments, and know when a comparator is legal;
  \item compress coordinates;
  \item prove greedy algorithms with ``greedy stays ahead'' and exchange arguments, and break wrong ones with
        counterexamples;
  \item combine greedy choices with heaps and ordered sets, including greedy algorithms that can undo their choices.
\end{chapterintro}

% ======================================================================================================================
\section[Sorting as a tool]{Sorting as a tool \hfill \bronze}\label{sec:sorting-tool}

\cc{sort} runs in $\Oh{n\log n}$: about 40\,ms for $10^6$ numbers (Table~\ref{tab:sortcost}), so for typical limits it is
almost free. What makes it so useful is what the order gives you afterwards.

\begin{itemize}
  \item \textbf{Equal values are adjacent.} Counting duplicates, grouping items by key, and counting distinct values all
        become one scan: after sorting, a new group starts exactly where \cc{a[i] != a[i - 1]}. The idiom
        \cc{a.erase(unique(all(a)), a.end())} removes duplicates.
  \item \textbf{Order statistics are positions.} The $k$-th smallest value is \cc{a[k - 1]}, the median is in the middle,
        and ``how many values are below $x$'' is a binary search (Chapter~\ref{ch:binsearch}).
  \item \textbf{Closest values are neighbours.} For any $i < k < j$ in sorted order, $a_j - a_i \ge a_k - a_i$, so the
        minimum difference between two values is attained by two adjacent values. The same holds for ``the value nearest to
        $x$'': it is next to where $x$ would be inserted.
  \item \textbf{Scans with two pointers become possible.} Once both of two lists are sorted, matching them (pairs with a given
        sum, assigning applicants to apartments) takes one linear pass (Chapter~\ref{ch:prefix}).
  \item \textbf{Sorting fixes the order of decisions.} A greedy algorithm needs to know which item to consider first; sorting
        by the right key is usually the whole difficulty.
\end{itemize}

A beautiful example of the last point, and of an invariant that holds throughout a scan of sorted data, is the following
problem.

\begin{problem}{Missing Coin Sum (CSES)}
Given $n \le 2\cdot10^5$ coins with positive values, find the smallest sum that cannot be formed by choosing a subset of the
coins.
\solution Sort the coins and maintain the invariant \emph{every sum $0, 1, \ldots, r$ can be formed with the coins seen
so far}; initially $r = 0$. Take the next coin $c$. If $c \le r + 1$, then adding $c$ to the sums $0, \ldots, r$ gives all
sums $c, \ldots, r + c$, which together with $0, \ldots, r$ cover $0, \ldots, r + c$ without a gap, so $r$ becomes $r + c$.
If $c > r + 1$, then $r + 1$ can never be formed: the coins seen so far sum to exactly $r$ (by induction, $r$ is the sum of
all coins so far), and every remaining coin is at least $c > r + 1$ by the sorted order. So the answer is $r + 1$
(Figure~\ref{fig:missingcoin}).
\end{problem}

\begin{figure}[htbp]
\centering
\begin{tikzpicture}[x=0.52cm, y=0.62cm]
  \draw[->, line width=0.6pt, draw=fgink!70] (-0.3, 0) -- (21.6, 0) node[right, lbl] {sum};
  \foreach \t in {0, 5, 10, 15, 20} \draw[line width=0.5pt, draw=fgink!70] (\t, 0.1) -- (\t, -0.1)
    node[below, lbl] {\t};
  \foreach \cl/\rr/\yy in {start/0/5, {after coin 1}/1/4, {after coin 1}/2/3, {after coin 3}/5/2,
      {after coin 6}/11/1} {
    \draw[ivl=fggreen!60] (0, \yy) -- (\rr + 0.02, \yy);
    \node[lbl, anchor=east] at (-0.5, \yy) {\cl: \ $r = \rr$};
  }
  \draw[line width=1pt, draw=fgverm, dashed] (12, -0.15) -- (12, 5.4);
  \node[font=\sffamily\scriptsize\bfseries, text=fgverm, anchor=south] at (12, 5.4) {12 cannot be formed};
  \fill[fgverm] (20, 1) circle (3pt) node[above=2pt, font=\sffamily\scriptsize, text=fgverm] {next coin: $20 > r + 1$};
\end{tikzpicture}
\caption{Missing Coin Sum with the sorted coins $1, 1, 3, 6, 20$. Each green bar is the range $0, \ldots, r$ of sums that can
be formed after taking the coins so far. Each coin at most $r + 1$ extends the range without a gap; the coin 20 would leave
the gap $12, \ldots, 19$, so the answer is 12.}
\label{fig:missingcoin}
\end{figure}

\begin{code}
void solve() {
    int n;
    cin >> n;
    vi x(n);
    cin >> x;
    sort(all(x));
    ll reach = 0;  // invariant: every sum 0, 1, ..., reach can be formed
    for (ll c : x) {
        if (c > reach + 1) break;  // then reach + 1 can never be formed
        reach += c;
    }
    print(reach + 1);
}
\end{code}

\subsection{The sorting toolbox}

The standard library offers more than \cc{sort}, and choosing the right routine matters when the input is large or only part
of the order is needed (Figure~\ref{fig:sortfamily}).

\begin{itemize}
  \item \cc{sort(all(a))} sorts in $\Oh{n \log n}$ (introsort: quicksort that falls back to heapsort, so the worst case is
        guaranteed). Equal elements may be reordered.
  \item \cc{stable_sort} keeps equal elements in their original order, at $\Oh{n\log n}$ with extra memory. It is what you
        need when sorting by one key after another (``by age, and among equal ages by the original order'').
  \item \cc{partial_sort(a.begin(), a.begin() + k, a.end())} puts the $k$ smallest elements, sorted, at the front, in
        $\Oh{n\log k}$. The rest is left in unspecified order.
  \item \cc{nth_element(a.begin(), a.begin() + k, a.end())} puts the element that belongs at position $k$ there, with smaller
        or equal elements before it and larger or equal ones after it, in $\Oh{n}$ on average: the fastest way to find a
        median or a threshold.
  \item Pairs, tuples and \cc{array}s compare lexicographically, so \cc{sort} on a vector of \cc{tuple<int, int, int>} sorts
        by the first component, then the second, then the third. To sort by the second component first, store it first.
  \item To sort indices by a key without moving the data (an \defn{argsort}), sort \cc{iota(all(ord), 0)} with a comparator
        that compares \cc{key[i]} and \cc{key[j]}. Keep the original indices whenever the output must refer to them.
\end{itemize}

\begin{figure}[htbp]
\centering
\begin{tikzpicture}
  \node[blbl, anchor=east] at (-0.25, 0) {input};
  \arrs{I}{0}{0}{7/n, 2/n, 9/n, 4/n, 1/n, 8/n, 3/n, 6/n}
  \node[blbl, anchor=east] at (-0.25, -1.05) {sort};
  \arrs{S}{0}{-1.05}{1/good, 2/good, 3/good, 4/good, 6/good, 7/good, 8/good, 9/good}
  \node[blbl, anchor=east] at (-0.25, -2.2) {partial\_sort, $k = 3$};
  \arrs{P}{0}{-2.2}{1/good, 2/good, 3/good, 9/dim, 7/dim, 8/dim, 4/dim, 6/dim}
  \node[blbl, anchor=east] at (-0.25, -3.5) {nth\_element, $k = 3$};
  \arrs{N}{0}{-3.5}{3/cool, 1/cool, 2/cool, 4/hot, 8/warn, 9/warn, 7/warn, 6/warn}
  \idx{N}{8}
  \node[notel, anchor=west, text width=6.6cm] at (6.0, -1.05) {all sorted, $\Oh{n\log n}$};
  \node[notel, anchor=west, text width=6.6cm] at (6.0, -2.2) {the 3 smallest, sorted, then the rest in some order;
    $\Oh{n\log k}$};
  \node[notel, anchor=west, text width=6.6cm] at (6.0, -3.5) {the right value at index 3; smaller values before it,
    larger ones after it, each side in some order; $\Oh{n}$};
\end{tikzpicture}
\caption{What \texttt{sort}, \texttt{partial\_sort} and \texttt{nth\_element} guarantee. Grey cells are in an unspecified
order; with \texttt{nth\_element}, blue cells are at most the value at index $k$ and orange cells at least it.}
\label{fig:sortfamily}
\end{figure}

\usetemplate{Argsort + custom comparator}{Ctrl+Alt+T, Shift+J}{cpargsort}

% ======================================================================================================================
\section[Counting sort and radix sort]{Counting sort and radix sort \hfill \silver}\label{sec:radix}

Comparison sorts cannot beat $\Oh{n \log n}$: there are $n!$ possible orders and each comparison has two outcomes, so some
input needs $\log_2 n! \approx n \log_2 n$ comparisons. Sorting \emph{integers} escapes this bound by not comparing at all.

\paragraph{Counting sort.} If the keys are small integers in $[0, K)$, count how many items have each key, turn the counts
into starting positions with a prefix sum, and place every item at the next free position of its key
(Figure~\ref{fig:countingsort}). This takes $\Oh{n + K}$ and is \emph{stable}: items with equal keys keep their input order,
because they are placed in input order. Stability is what makes the next algorithm work.

\begin{figure}[htbp]
\centering
\begin{tikzpicture}[kc/.style={cellc, minimum width=8mm}]
  \tikzset{k0/.style={fill=fgsky!30, draw=fgblue}, k1/.style={fill=fggreen!22, draw=fggreen},
    k2/.style={fill=fgorange!25, draw=fgorange!85!black}, k3/.style={fill=fgpink!25, draw=fgpink!80!black}}
  \node[blbl, anchor=east] at (-0.3, 0) {input (key, item)};
  \foreach \kv/\lt/\s [count=\i] in {2/a/k2, 0/b/k0, 3/c/k3, 2/d/k2, 0/e/k0, 1/f/k1, 2/g/k2}
    \node[kc, \s] (in\i) at (\i*0.8 - 0.4, 0) {\kv\,\textit{\lt}};
  \node[blbl, anchor=east] at (-0.3, -1.3) {count per key};
  \foreach \c/\s [count=\i] in {2/k0, 1/k1, 3/k2, 1/k3} {
    \pgfmathtruncatemacro{\kk}{\i - 1}
    \node[kc, \s] (cnt\i) at (\i*0.8 - 0.4, -1.3) {\c};
    \node[lbl, above=0pt of cnt\i] {key \kk};
  }
  \node[blbl, anchor=east] at (-0.3, -2.45) {first position};
  \foreach \c/\s [count=\i] in {0/k0, 2/k1, 3/k2, 6/k3} \node[kc, \s] at (\i*0.8 - 0.4, -2.45) {\c};
  \node[notel, anchor=west, text width=5.2cm] at (3.6, -1.9) {prefix sums of the counts: key 2 starts at position
    $2 + 1 = 3$};
  \node[blbl, anchor=east] at (-0.3, -3.75) {output};
  \foreach \kv/\lt/\s [count=\i] in {0/b/k0, 0/e/k0, 1/f/k1, 2/a/k2, 2/d/k2, 2/g/k2, 3/c/k3}
    \node[kc, \s] (out\i) at (\i*0.8 - 0.4, -3.75) {\kv\,\textit{\lt}};
  \foreach \i in {0,...,6} \node[lbl] at (\i*0.8 + 0.4, -4.3) {\i};
  \node[notel, anchor=west, text width=5.2cm] at (6.1, -3.75) {items with equal keys keep their input order
    ($a, d, g$): the sort is stable};
\end{tikzpicture}
\caption{Counting sort with keys in $[0, 4)$. Count the keys, turn the counts into first positions with a prefix sum, then
place the items in input order, each at the next free position of its key.}
\label{fig:countingsort}
\end{figure}

\paragraph{Radix sort.} To sort $d$-digit numbers, counting-sort them by the \emph{last} digit, then by the second-to-last,
and so on up to the first digit (Figure~\ref{fig:radix}). This is \defn{radix sort}, least significant digit first. It is
correct by induction: after the pass on digit $j$, the numbers are sorted by their last $j$ digits, because the pass orders
them by digit $j$, and numbers that tie on digit $j$ keep their previous order, which by induction was sorted by the last
$j - 1$ digits. Each pass costs $\Oh{n + b}$ for base $b$.

\begin{figure}[htbp]
\centering
\begin{tikzpicture}[dg/.style={font=\ttfamily\small, inner sep=0.8pt}]
  \foreach \col/\vals/\hl/\ttl in {0/{329,457,657,839,436,720,355}/-1/input,
      1/{720,355,436,457,657,329,839}/2/{by ones},
      2/{720,329,436,839,355,457,657}/1/{by tens},
      3/{329,355,436,457,657,720,839}/0/{by hundreds}} {
    \node[title] at (\col*2.9, 0.55) {\ttl};
    \foreach \vv [count=\rw] in \vals {
      \node[draw=fgline, line width=0.5pt, fill=white, minimum width=13mm, minimum height=5.2mm]
        at (\col*2.9, -\rw*0.58) {};
      \pgfmathtruncatemacro{\da}{floor(\vv/100)}
      \pgfmathtruncatemacro{\db}{mod(floor(\vv/10),10)}
      \pgfmathtruncatemacro{\dd}{mod(\vv,10)}
      \ifnum\hl=0 \node[dg, text=fgverm, font=\ttfamily\small\bfseries] at (\col*2.9-0.2, -\rw*0.58) {\da};
      \else \node[dg] at (\col*2.9-0.2, -\rw*0.58) {\da}; \fi
      \ifnum\hl=1 \node[dg, text=fgverm, font=\ttfamily\small\bfseries] at (\col*2.9, -\rw*0.58) {\db};
      \else \node[dg] at (\col*2.9, -\rw*0.58) {\db}; \fi
      \ifnum\hl=2 \node[dg, text=fgverm, font=\ttfamily\small\bfseries] at (\col*2.9+0.2, -\rw*0.58) {\dd};
      \else \node[dg] at (\col*2.9+0.2, -\rw*0.58) {\dd}; \fi
    }
  }
  \foreach \a/\b in {0/1, 1/2, 2/3} \draw[farr] (\a*2.9+0.85, -2.32) -- (\b*2.9-0.85, -2.32);
\end{tikzpicture}
\caption{Radix sort of three-digit numbers: three stable counting sorts, by the ones digit, the tens digit and the hundreds
digit (the digit used in each pass is red). After the pass on the tens digit, the numbers are sorted by their last two
digits; after the last pass, completely.}
\label{fig:radix}
\end{figure}

For 64-bit integers the template uses base $2^{16}$, four passes of 16 bits each, and flips the sign bit first so that
negative numbers sort correctly as unsigned values. Table~\ref{tab:sortcost} shows that it is more than twice as fast as
\cc{sort} on large inputs. That rarely decides a contest problem, but it does when $n$ is $10^7$, when the same data is
sorted many times, or inside a suffix array construction (Chapter~\ref{ch:strings}), which is a radix sort of pairs of ranks.

\begin{table}[htbp]
\centering
\small
\begin{tabular}{@{}l r r l@{}}
  \toprule
  \thd{Routine} & \thd{$n = 10^6$} & \thd{$n = 10^7$} & \thd{Guarantee} \\
  \midrule
  \texttt{sort} & 43\,ms & 515\,ms & $\Oh{n\log n}$ \\
  \texttt{stable\_sort} & 51\,ms & 615\,ms & $\Oh{n\log n}$, stable \\
  radix sort (template) & 17\,ms & 238\,ms & $\Oh{n}$, 4 passes, stable \\
  \texttt{nth\_element} (median) & 6\,ms & 48\,ms & $\Oh{n}$ on average \\
  \texttt{partial\_sort}, $k = 100$ & $< 1$\,ms & 2\,ms & $\Oh{n \log k}$ \\
  \bottomrule
\end{tabular}
\caption{Measured costs of sorting random 64-bit integers (\texttt{g++ -O2}, best of three runs).}
\label{tab:sortcost}
\end{table}

\usetemplate{Radix sort and counting sort}{Ctrl+Alt+U, Shift+R}{cpradix}

% ======================================================================================================================
\section[Comparators and exchange arguments]{Comparators and exchange arguments \hfill \silver}\label{sec:comparators}

When the order you need is not the natural order of the values, you write it as a comparator: a function \cc{cmp(a, b)} that
returns \cc{true} when $a$ must come before $b$. Finding the comparator is the creative part, and the tool for it is the
\emph{adjacent exchange argument} of Section~\ref{sec:exchange}: take an optimal order, look at two neighbours $x$ and $y$,
and ask when swapping them helps. Because only the two neighbours are affected, the condition usually becomes a simple
inequality in the parameters of $x$ and $y$, and that inequality \emph{is} the comparator.

\begin{problem}{Weighted completion time}
Task $i$ takes $t_i$ time units and has weight $w_i$ (both positive). The tasks run one after another, and the cost is
$\sum_i w_i C_i$, where $C_i$ is the time at which task $i$ finishes. Find the minimum cost.
\solution Consider two adjacent tasks $x$ and $y$ that start at time $T$ (Figure~\ref{fig:smith}). Swapping them does not
change when any other task finishes, so only their own two terms change. With $x$ first, the two terms are
$w_x (T + t_x) + w_y (T + t_x + t_y)$; with $y$ first, $w_y (T + t_y) + w_x (T + t_x + t_y)$. The difference is
$w_y t_x - w_x t_y$, so $x$ should come first exactly when $t_x w_y < t_y w_x$, that is, when $t_x / w_x < t_y / w_y$. Sort
by this ratio (\defn{Smith's rule}); any order that is not sorted has an adjacent pair that can be swapped without increasing
the cost, so the sorted order is optimal.
\end{problem}

\begin{figure}[htbp]
\centering
\begin{tikzpicture}[x=0.55cm, y=0.8cm]
  \tikzset{jb/.style={draw=fgink!60, line width=0.6pt, minimum height=6mm, inner sep=0pt, font=\sffamily\small}}
  \foreach \yy/\lab/\wa/\wb/\na/\nb/\ca/\cb in {0/{$x$ first}/4/6/x/y/fgblue/fgorange,
      -1.5/{$y$ first}/6/4/y/x/fgorange/fgblue} {
    \node[blbl, anchor=east] at (-0.3, \yy) {\lab};
    \node[jb, fill=fgline!30, minimum width=5*0.55cm, anchor=west] at (0, \yy) {earlier tasks};
    \node[jb, fill=\ca!25, minimum width=\wa*0.55cm, anchor=west] at (5, \yy) {$\na$};
    \node[jb, fill=\cb!25, minimum width=\wb*0.55cm, anchor=west] at (5 + \wa, \yy) {$\nb$};
    \node[jb, fill=fgline!30, minimum width=5*0.55cm, anchor=west] at (15, \yy) {later tasks};
  }
  \draw[line width=0.6pt, draw=fgink!60, dashed] (5, 0.55) -- (5, -2.05) node[below, lbl] {$T$};
  \draw[line width=0.6pt, draw=fgink!60, dashed] (15, 0.55) -- (15, -2.05) node[below, lbl] {$T + t_x + t_y$};
  \node[notel, anchor=north, text width=10cm, align=center] at (10, -2.75) {Only $x$ and $y$ finish at different times,
    so the swap changes the cost by $w_y t_x - w_x t_y$.};
\end{tikzpicture}
\caption{The exchange argument behind Smith's rule. Swapping two adjacent tasks changes only their own finishing times; the
tasks before and after them are unaffected.}
\label{fig:smith}
\end{figure}

\begin{code}
void solve() {
    int n;
    cin >> n;
    vector<pll> job(n);  // (t, w): duration and weight, both positive
    for (auto& [t, w] : job) cin >> t >> w;
    // x goes first exactly when t_x * w_y < t_y * w_x (Smith's rule)
    sort(all(job), [](const pll& x, const pll& y) {
        return x.first * y.second < y.first * x.second;
    });
    ll time = 0, cost = 0;
    for (auto [t, w] : job) {
        time += t;  // this job finishes at the current time
        cost += w * time;
    }
    print(cost);
}
\end{code}

The comparator cross-multiplies instead of comparing the ratios $t/w$ as doubles: integer arithmetic is exact, and with
$t, w \le 10^6$ the products fit comfortably in 64 bits. With all weights equal, Smith's rule becomes \emph{shortest job
first}, which minimises the sum (and so the average) of completion times.

\subsection{When is a comparator legal?}

\cc{sort} assumes that the comparator is a \defn{strict weak ordering}: $\mathrm{cmp}(a, a)$ is false; if $\mathrm{cmp}(a,
b)$ and $\mathrm{cmp}(b, c)$ then $\mathrm{cmp}(a, c)$; and ``neither comes before the other'' is transitive too, so that
equivalent elements form groups. If any of these fails, the behaviour is undefined, and in practice \cc{sort} can produce an
unsorted result, loop forever or read outside the array. An exchange argument only tells you how to order \emph{two}
elements; before trusting it as a comparator, check that it is transitive, usually by rewriting it as a comparison of a single
key computed from each element, like $t/w$ above.

\begin{problem}{The largest concatenation}
Arrange $n$ non-negative integers in a row so that their concatenation, read as one number, is as large as possible (for
$5, 50, 56$ the answer is $56550$).
\solution Adjacent exchange: $a$ should precede $b$ when the string $a + b$ is larger than $b + a$. This is transitive,
because $a + b > b + a$ as numbers means $a \cdot 10^{|b|} + b > b \cdot 10^{|a|} + a$, which rearranges to
$a / (10^{|a|} - 1) > b / (10^{|b|} - 1)$: a comparison of one key per number. So sorting with it is legal and, by the
exchange argument, optimal. If the largest element is 0, all are, and the answer is the single digit 0.
\end{problem}

\begin{code}
void solve() {
    int n;
    cin >> n;
    vector<string> s(n);
    cin >> s;
    // a goes first when the concatenation a + b beats b + a
    sort(all(s), [](const string& a, const string& b) {
        return a + b > b + a;
    });
    if (s[0] == "0") s = {"0"};  // every number is zero
    string out;
    for (auto& x : s) out += x;
    print(out);
}
\end{code}

\begin{pitfall}[Comparators that break \texttt{sort}]
Never use \cc{<=} in a comparator: $\mathrm{cmp}(a, a)$ becomes true, and GCC's \cc{sort} can then run past the end of the
array. Never compare floating-point keys that may be NaN. And be careful with comparators that look symmetric but are not
transitive, such as ``$a$ before $b$ if $a$ beats $b$'' in a rock-paper-scissors tournament: no order is consistent with
them. When in doubt, compile with \cc{-D_GLIBCXX_DEBUG}, which checks the comparator on every call.
\end{pitfall}

% ======================================================================================================================
\section[Coordinate compression]{Coordinate compression \hfill \silver}\label{sec:compression}

Many data structures index arrays by value: counting arrays, Fenwick trees, segment trees. When values are huge (up to $10^9$)
but there are few of them ($2\cdot10^5$), \defn{coordinate compression} replaces every value by its \emph{rank} among the
distinct values (Figure~\ref{fig:compress}). The order of values is preserved, equal values stay equal, and the new values
are small enough to index arrays.

\begin{figure}[htbp]
\centering
\begin{tikzpicture}
  \node[blbl, anchor=east] at (-0.25, 0) {values};
  \arrs{A}{0}{0}{1e9/n, 7/n, 42/n, 7/n, 999/n}
  \node[blbl, anchor=east] at (-0.25, -1.6) {ranks};
  \arrs{B}{0}{-1.6}{3/cool, 0/cool, 1/cool, 0/cool, 2/cool}
  \foreach \i in {1,...,5} \draw[->, line width=0.8pt, draw=fgblue] (A-\i.south) -- (B-\i.north);
  \node[blbl, anchor=east] at (6.4, -0.8) {sorted distinct};
  \arrs{V}{6.6}{-0.8}{7/n, 42/n, 999/n, 1e9/n}
  \idx{V}{4}
  \node[notel, anchor=north west, text width=4.4cm] at (6.6, -1.75) {rank of $v$ = position of $v$ here, found with
    \texttt{lower\_bound}};
\end{tikzpicture}
\caption{Coordinate compression: every value is replaced by its index in the sorted list of distinct values.}
\label{fig:compress}
\end{figure}

In code: copy the values, sort, remove duplicates, and map each value $v$ to \cc{lower_bound(all(vals), v) - vals.begin()}.
Building costs $\Oh{n\log n}$ and each lookup $\Oh{\log n}$; if lookups are frequent, store the ranks once in an array.

\usetemplate{Coordinate compression}{Ctrl+Alt+T, Shift+C}{cpcompress}

\begin{pitfall}[Compress everything that will be looked up]
Offline problems often have values in the data \emph{and} in the queries. Collect all of them before compressing, or a query
value that does not occur in the data has no rank. When intervals are compressed and the gaps between endpoints matter
(``is the point between two endpoints covered?''), also insert $v + 1$ for each endpoint $v$, or keep the original values next
to the ranks to measure lengths.
\end{pitfall}

% ======================================================================================================================
\section[Greedy algorithms and their proofs]{Greedy algorithms and their proofs \hfill \silver}\label{sec:greedy}

A \defn{greedy algorithm} builds a solution step by step, always taking the choice that looks best now and never
reconsidering it. Greedy algorithms are short and fast, and they are wrong more often than they are right: many natural greedy
rules fail on small inputs. So a greedy algorithm needs either a proof or, at the very least, a stress test against a brute
force (Chapter~\ref{ch:approach}). Two proof patterns cover almost every contest greedy.

\begin{itemize}
  \item \textbf{Greedy stays ahead.} Show that after each step, the greedy partial solution is at least as good as the
        corresponding part of \emph{any} solution, by some measure; then it is also at least as good at the end.
  \item \textbf{Exchange argument.} Take an optimal solution that differs from the greedy one, and transform it step by step
        into the greedy solution, never making it worse (Section~\ref{sec:exchange}). Then the greedy solution is optimal
        too.
\end{itemize}

\subsection{Interval scheduling}

\begin{problem}{Most disjoint intervals}
Given $n$ intervals, choose as many as possible so that no two of them overlap.
\solution Sort the intervals by their end. Scan them, taking every interval that starts no earlier than the end of the last
interval taken (Figure~\ref{fig:intervalsched}). This takes $\Oh{n\log n}$; the CSES problem Movie Festival is exactly this
(Chapter~\ref{ch:worked}).
\end{problem}

\emph{Why it works (greedy stays ahead).} Let $g_1, g_2, \ldots, g_k$ be the intervals chosen by the greedy algorithm and
$o_1, o_2, \ldots, o_m$ those of any valid solution, both in order of time. We claim that $g_i$ ends no later than $o_i$ for
every $i \le k$. For $i = 1$ this holds because $g_1$ has the earliest end of all intervals. If $g_{i}$ ends no later than
$o_i$, then $o_{i+1}$, which starts after $o_i$ ends, also starts after $g_i$ ends, so $o_{i+1}$ was still available when
the greedy algorithm chose $g_{i+1}$, and the greedy algorithm chose the available interval with the earliest end: $g_{i+1}$
ends no later than $o_{i+1}$. Finally, if $m > k$, the interval $o_{k+1}$ starts after $o_k$ ends, hence after $g_k$ ends,
so the greedy algorithm would have taken another interval: a contradiction. So $k \ge m$.

\begin{figure}[htbp]
\centering
\begin{tikzpicture}[x=0.95cm, y=0.44cm]
  \numline{0}{12}{0}{time}
  \foreach \l/\r/\y/\c in {0/3/1/fggreen!70, 2/4/2/fgline!70, 1/5/3/fgline!70, 3.5/6/4/fggreen!70, 6.5/8/5/fggreen!70,
      4/9/6/fgline!70, 7/11/7/fgline!70, 8.5/12/8/fggreen!70}
    \draw[ivl=\c] (\l, \y) -- (\r, \y);
  \foreach \e/\y in {3/1, 6/4, 8/5, 12/8} \draw[line width=0.6pt, draw=fggreen!60!black, dashed] (\e, \y) -- (\e, 0.2);
  \node[notel, anchor=west, text width=3.3cm] at (12.8, 4.5) {one row per interval, by end time from bottom to top;
    green: taken};
\end{tikzpicture}
\caption{Earliest end first. Each interval is on its own row, in order of end time. The algorithm takes an interval whenever
it starts at or after the end of the last one taken (dashed lines): four intervals, the maximum.}
\label{fig:intervalsched}
\end{figure}

\usetemplate{Intervals: merge + max disjoint (greedy)}{Ctrl+Alt+T, Shift+N}{cpintervals}

The same greedy, with the same proof, solves the \emph{stabbing} problem: the fewest points such that every interval contains
at least one of them. Sort by end, and whenever an interval does not contain the last point placed, place a new point at its
end. The number of points equals the maximum number of disjoint intervals, an instance of a min-max duality that reappears
with matchings and flows (Chapter~\ref{ch:flows}).

\subsection{Breaking a wrong greedy}

Plausible greedy rules fail all the time, and the fastest way to discover it is to look for a small counterexample, by hand or
with a stress test. Figure~\ref{fig:greedyfail} shows three classics. The lesson is not that greedy is bad, but that the
\emph{specific} rule matters, and that a rule is only as good as its proof.

\begin{figure}[htbp]
\centering
\begin{tikzpicture}[x=0.36cm, y=0.5cm]
  \begin{scope}
    \node[title, anchor=west] at (0, 3.1) {(a) shortest first};
    \draw[ivl=fgline!70] (0, 1.9) -- (5, 1.9);
    \draw[ivl=fgline!70] (6, 1.9) -- (11, 1.9);
    \draw[ivl=fgverm!70] (4, 0.9) -- (7, 0.9);
    \node[notel, anchor=north west, text width=3.9cm] at (0, 0.2) {Greedy takes the short red interval and then nothing
      else; the two grey ones are better.};
  \end{scope}
  \begin{scope}[xshift=5.4cm]
    \node[title, anchor=west] at (0, 3.1) {(b) earliest start first};
    \draw[ivl=fgverm!70] (0, 1.9) -- (11, 1.9);
    \draw[ivl=fgline!70] (0.5, 0.9) -- (3, 0.9);
    \draw[ivl=fgline!70] (4, 0.9) -- (6.5, 0.9);
    \draw[ivl=fgline!70] (7.5, 0.9) -- (10, 0.9);
    \node[notel, anchor=north west, text width=3.9cm] at (0, 0.2) {Greedy takes the long interval that starts first; the
      three short ones are better.};
  \end{scope}
  \begin{scope}[xshift=10.8cm]
    \node[title, anchor=west] at (0, 3.1) {(c) largest coin first};
    \foreach \cv/\cx in {4/0.9, 1/3.1, 1/5.1}
      \node[circle, draw=fgverm, fill=fgverm!15, line width=0.8pt, minimum size=6mm, inner sep=0pt,
        font=\sffamily\small] at (\cx, 1.9) {\cv};
    \node[lbl, anchor=west] at (6.4, 1.9) {greedy: 3};
    \foreach \cv/\cx in {3/0.9, 3/3.1}
      \node[circle, draw=fggreen, fill=fggreen!15, line width=0.8pt, minimum size=6mm, inner sep=0pt,
        font=\sffamily\small] at (\cx, 0.6) {\cv};
    \node[lbl, anchor=west] at (6.4, 0.6) {best: 2};
    \node[notel, anchor=north west, text width=3.9cm] at (0, -0.3) {Coins $1, 3, 4$ and the amount 6.};
  \end{scope}
\end{tikzpicture}
\caption{Three plausible greedy rules and their counterexamples. Interval scheduling needs the earliest \emph{end}; coin
change needs dynamic programming unless the coin system is special (such as $1, 2, 5, 10, \ldots$).}
\label{fig:greedyfail}
\end{figure}

% ======================================================================================================================
\section[Classic greedy problems]{Classic greedy problems \hfill \silver}\label{sec:greedy-classics}

A small number of greedy ideas account for a large share of contest problems. This section works through two of them in
detail, each with its proof, because the proofs transfer: the slope argument for medians reappears in slope trick
(Chapter~\ref{ch:dpopt}), and the exchange argument behind Huffman's algorithm is the template for every ``merge the two
smallest'' problem. A catalogue of the other classic rules closes the section.

\subsection{Medians minimise distances}

\begin{problem}{Stick Lengths (CSES)}
Given $n$ sticks with lengths $p_1, \ldots, p_n$, make all sticks equally long; changing a length by one costs one. Find the
minimum total cost.
\solution The cost of the target length $x$ is $f(x) = \sum_i |p_i - x|$. Moving $x$ one unit to the right increases the
distance to every stick at or left of $x$ and decreases it for every stick right of $x$, so the slope of $f$ is (number of
sticks to the left) $-$ (number to the right). This slope is negative while fewer than half of the sticks are to the left and
positive afterwards, so $f$ is minimised at a \emph{median} (Figure~\ref{fig:sticks}). Any median works, and
\cc{nth_element} finds one in linear time.
\end{problem}

\begin{figure}[htbp]
\centering
\begin{tikzpicture}[slope/.style={font=\sffamily\scriptsize, text=fgink!80, fill=white, inner sep=1pt,
    anchor=south east}]
  \begin{axis}[cpplot, width=10cm, height=5.2cm, xmin=-0.3, xmax=10.3, ymin=8, ymax=30, xlabel={target length $x$},
      ylabel={total cost $f(x)$}, xtick={0,...,10}, ytick={10,15,20,25,30}]
    \addplot[fgblue, line width=1.3pt, mark=*, mark size=1.6pt] coordinates {(0,22) (1,17) (2,14) (3,13) (4,14) (5,15)
      (6,16) (7,17) (8,20) (9,23) (10,28)};
    \addplot[only marks, mark=*, mark size=3pt, fgverm] coordinates {(3,13)};
    \node[font=\sffamily\scriptsize, text=fgverm, anchor=north west] at (axis cs:3.2,12.3) {median 3: cost 13};
    \node[slope, anchor=north east] at (axis cs:0.4,18.7) {$-5$};
    \node[slope, anchor=north east] at (axis cs:1.4,14.7) {$-3$};
    \node[slope, anchor=north east] at (axis cs:2.4,12.7) {$-1$};
    \node[slope, anchor=north west] at (axis cs:5.1,14.2) {$+1$};
    \node[slope, anchor=north west] at (axis cs:8.1,19.2) {$+3$};
    \node[slope, anchor=north west] at (axis cs:9.6,24.7) {$+5$};
  \end{axis}
\end{tikzpicture}
\caption{The cost $f(x) = \sum |p_i - x|$ for sticks $1, 2, 3, 7, 9$. It is piecewise linear and convex; its slope (numbers
on the segments) goes up by 2 at each stick, and changes sign at the median.}
\label{fig:sticks}
\end{figure}

\begin{code}
void solve() {
    int n;
    cin >> n;
    vi p(n);
    cin >> p;
    // a median minimises the sum of distances
    nth_element(p.begin(), p.begin() + n / 2, p.end());
    ll med = p[n / 2], cost = 0;
    for (ll x : p) cost += abs(x - med);
    print(cost);
}
\end{code}

The same argument, on each coordinate separately, shows that the point minimising the sum of Manhattan distances to a set of
points in the plane is the coordinate-wise median. With squared distances instead, the optimum is the mean.

\subsection{Merging with the smallest first: Huffman's algorithm}

\begin{problem}{Stick Divisions (CSES)}
A stick of length $x$ must be cut into $n$ pieces of given lengths $d_1, \ldots, d_n$ (which sum to $x$). Cutting a stick of
length $\ell$ into two parts costs $\ell$. Find the minimum total cost.
\solution Run the process backwards: start from the pieces and \emph{join} two sticks at a time, paying the length of the
result. The sequence of cuts is a binary tree whose leaves are the pieces and whose internal nodes are the cut sticks, and the
total cost is the sum of the internal nodes, which equals $\sum_i d_i \cdot \mathrm{depth}_i$ (Figure~\ref{fig:huffman}).
This is exactly the cost of a prefix code with frequencies $d_i$, and \defn{Huffman's algorithm} minimises it: repeatedly
join the two shortest sticks.
\end{problem}

\emph{Why joining the two shortest is safe (exchange argument).} In an optimal tree, take a deepest internal node; its two
children are leaves at the maximum depth. Swapping these leaves with the two shortest pieces cannot increase
$\sum d_i\,\mathrm{depth}_i$, since the shortest pieces move deeper and the longer ones move up. So some optimal tree has the
two shortest pieces as siblings; joining them first and treating the result as a single piece of length $d_a + d_b$ leaves a
smaller instance of the same problem.

\begin{figure}[htbp]
\centering
\begin{forest}
  for tree={rectangle, rounded corners=2pt, draw=fgblue, fill=fgfill, line width=0.7pt, font=\sffamily\small,
    minimum width=8mm, minimum height=6mm, inner sep=2pt, s sep=6mm, l sep=6mm, edge={line width=0.8pt,
    draw=fgline!80!fgink}}
  [18, cool, label={right:{\scriptsize\sffamily\color{fgblue} cut, cost 18}}
    [8, warn]
    [10, cool, label={right:{\scriptsize\sffamily\color{fgblue} cost 10}}
      [5, cool, label={left:{\scriptsize\sffamily\color{fgblue} cost 5}}
        [2, warn] [3, warn]]
      [5, warn]]]
\end{forest}
\caption{Cutting a stick of length 18 into pieces $2, 3, 5, 8$ (orange). Huffman's algorithm joins $2 + 3 = 5$, then
$5 + 5 = 10$, then $10 + 8 = 18$; the cuts (blue) cost $5 + 10 + 18 = 33 = 8 \cdot 1 + 5 \cdot 2 + 2 \cdot 3 + 3 \cdot 3$.}
\label{fig:huffman}
\end{figure}

\begin{code}
void solve() {
    int x, n;
    cin >> x >> n;
    minpq<ll> pq;
    rep(i, n) {
        ll d;
        cin >> d;
        pq.push(d);
    }
    ll cost = 0;
    while (sz(pq) > 1) {  // undo the cuts: always join the two shortest
        ll a = pq.top();
        pq.pop();
        ll b = pq.top();
        pq.pop();
        cost += a + b;  // cutting a + b into a and b costs a + b
        pq.push(a + b);
    }
    print(cost);
}
\end{code}

\subsection{A catalogue of classic greedy rules}

Table~\ref{tab:greedyrules} lists the greedy rules that every contestant should know by heart, together with the proof
technique. Each has been used in countless problems, often in disguise.

\begin{table}[htbp]
\centering
\small
\begin{tabular}{@{}P{5.1cm} P{6.2cm} P{3.4cm}@{}}
  \toprule
  \thd{Problem} & \thd{Greedy rule} & \thd{Proof} \\
  \midrule
  most disjoint intervals & earliest end first & stays ahead \\
  fewest points stabbing all intervals & sort by end; a point at the end of each unstabbed interval & stays ahead \\
  fewest rooms for all meetings & by start time; reuse the room that frees earliest (heap) & depth bound \\
  minimise the sum of completion times & shortest job first (Smith's rule with weights) & adjacent exchange \\
  minimise the maximum lateness & earliest deadline first & adjacent exchange \\
  merge piles, cost = size of result & join the two smallest (Huffman) & exchange \\
  equalise values, cost = distance & move everything to a median & slope argument \\
  fractional knapsack & by value per unit of weight & exchange \\
  coin change, canonical systems & largest coin that fits & special structure \\
  unit jobs with deadlines and profits & by deadline; keep the best in a heap & exchange (matroid) \\
  \bottomrule
\end{tabular}
\caption{Classic greedy rules. The 0/1 knapsack and coin change with arbitrary coins are \emph{not} solved by greedy rules;
they need dynamic programming (Part~V).}
\label{tab:greedyrules}
\end{table}

% ======================================================================================================================
\section[Greedy with heaps and ordered sets]{Greedy with heaps and ordered sets \hfill \gold}\label{sec:greedy-ds}

Many greedy algorithms process items in a fixed order (by time, by deadline) and at each step choose among a changing set of
candidates: the cheapest item so far, the member who became free last. A priority queue or an ordered multiset holds the
candidates, and the greedy choice becomes a heap operation or a \cc{lower_bound}.

\subsection{Deadlines and profits}

\begin{problem}{Scheduling unit jobs}
There are $n$ jobs, each taking one day, with deadline $d_i$ and profit $p_i$; a job earns its profit only if it is finished by
its deadline (on day $d_i$ or earlier), and one job can be done per day. Maximise the total profit.
\solution A set of jobs can be scheduled if and only if, for every $d$, at most $d$ of them have deadline at most $d$ (then
doing them in order of deadline works). Scan the jobs in order of deadline and keep the chosen profits in a min-heap. Add each
job; if the heap now holds more jobs than the current deadline allows, remove the smallest profit (Figure~\ref{fig:jobheap}).
The removed job is the one that is least useful, and every job still in the heap can be scheduled.
\end{problem}

\begin{figure}[htbp]
\centering
\small
\begin{tabular}{@{}c c l c l@{}}
  \toprule
  \thd{Job $(d, p)$} & \thd{Limit $d$} & \thd{Heap after adding} & \thd{Removed} & \thd{Chosen profits} \\
  \midrule
  $(1, 5)$ & 1 & $\{5\}$ & -- & $5$ \\
  $(2, 3)$ & 2 & $\{3, 5\}$ & -- & $3, 5$ \\
  $(2, 8)$ & 2 & $\{3, 5, 8\}$ & \textcolor{fgverm}{3} & $5, 8$ \\
  $(3, 1)$ & 3 & $\{1, 5, 8\}$ & -- & $1, 5, 8$ \\
  $(3, 6)$ & 3 & $\{1, 5, 6, 8\}$ & \textcolor{fgverm}{1} & $5, 6, 8$ \\
  $(3, 4)$ & 3 & $\{4, 5, 6, 8\}$ & \textcolor{fgverm}{4} & $5, 6, 8$ \quad (total 19) \\
  \bottomrule
\end{tabular}
\caption{Scheduling unit jobs with a min-heap. Whenever the heap holds more jobs than days available, the cheapest job is
dropped. The final jobs are done on days 1, 2, 3 in order of deadline: profits 5, 8, 6.}
\label{fig:jobheap}
\end{figure}

\begin{code}
void solve() {
    int n;
    cin >> n;
    vector<pll> job(n);  // (deadline, profit); each job takes one day
    for (auto& [d, p] : job) cin >> d >> p;
    sort(all(job));
    minpq<ll> kept;  // profits of the jobs chosen so far
    for (auto [d, p] : job) {
        kept.push(p);
        if (sz(kept) > d) kept.pop();  // only d jobs fit by day d
    }
    ll total = 0;
    for (; !kept.empty(); kept.pop()) total += kept.top();
    print(total);
}
\end{code}

This problem is the standard example of a \emph{matroid}: the sets of jobs that can be scheduled satisfy the exchange property
(a smaller schedulable set can always be extended by some job of a larger one), and on every matroid, taking elements in order
of decreasing weight whenever they keep the set independent is optimal. Kruskal's algorithm for minimum spanning trees
(Chapter~\ref{ch:dsu}) is the other famous instance.

\subsection{Best fit with an ordered multiset}

\begin{problem}{Movie Festival II (CSES)}
A club of $k$ members wants to watch as many of $n$ movies as possible (counting each watched movie once); each member can
watch any number of movies, one at a time, and a movie that ends at time $t$ does not overlap one that starts at time $t$.
\solution Process the movies by end time, as for a single person. For each movie, some member must be free at its start
time $a$; among the members who are, give the movie to the one who \emph{became free last} (the largest free time $\le a$).
This \emph{best fit} leaves the members with earlier free times available for movies that start earlier, and an exchange
argument shows that it never does worse than any other choice. An ordered multiset of free times finds that member with
\cc{upper_bound} in $\Oh{\log k}$.
\end{problem}

\begin{code}
void solve() {
    int n, k;
    cin >> n >> k;
    vector<pll> mv(n);  // (end, start), so that sorting orders by end
    for (auto& [b, a] : mv) cin >> a >> b;
    sort(all(mv));
    multiset<ll> freeAt;  // when each club member becomes free
    rep(i, k) freeAt.insert(0);
    int watched = 0;
    for (auto [b, a] : mv) {
        auto it = freeAt.upper_bound(a);  // first member free after a
        if (it == freeAt.begin()) continue;  // nobody is free at time a
        freeAt.erase(prev(it));  // the member who became free last
        freeAt.insert(b);
        watched++;
    }
    print(watched);
}
\end{code}

\subsection{Greedy that can change its mind}

Some problems defeat every fixed greedy rule because an early choice only turns out to be wrong much later. \defn{Regret
greedy} (also called greedy with undo) keeps that later option open: whenever it makes a choice, it also pushes a special
candidate into the heap that represents \emph{undoing} the choice, priced so that taking it later has exactly the effect of
having chosen differently.

\begin{problem}{Buy Low Sell High (Codeforces 865D)}
The price of a share is known for each of $n$ days. Each day you may buy one share, sell one share, or do nothing; you start
with no shares. Maximise the profit.
\solution Keep a min-heap of prices. On each day with price $p$, if the cheapest price $h$ in the heap is below $p$, sell a
share bought at $h$ for a profit of $p - h$, remove $h$, and push $p$ \emph{twice}. One copy is today as an ordinary buying
opportunity. The other copy is the regret option: if a later day has price $q > p$, ``buying at $p$ and selling at $q$''
earns $q - p$, and added to the $p - h$ already earned it gives $q - h$, exactly the profit of selling the share bought at $h$
on the later day instead. If $h \ge p$, just push $p$ (Figure~\ref{fig:regret}).
\end{problem}

\begin{figure}[htbp]
\centering
\small
\begin{tabular}{@{}c l l r@{}}
  \toprule
  \thd{Price} & \thd{Action} & \thd{Heap afterwards} & \thd{Profit} \\
  \midrule
  3 & heap empty: push 3 & $\{3\}$ & 0 \\
  1 & $3 \ge 1$: push 1 & $\{1, 3\}$ & 0 \\
  4 & sell the share bought at 1 ($+3$); push 4 twice & $\{3, 4, 4\}$ & 3 \\
  5 & sell the share bought at 3 ($+2$); push 5 twice & $\{4, 4, 5, 5\}$ & 5 \\
  2 & $4 \ge 2$: push 2 & $\{2, 4, 4, 5, 5\}$ & 5 \\
  6 & sell the share bought at 2 ($+4$); push 6 twice & $\{4, 4, 5, 5, 6, 6\}$ & 9 \\
  \bottomrule
\end{tabular}
\caption{Regret greedy on the prices $3, 1, 4, 5, 2, 6$. The optimal plan buys on days 1, 2 and 5 and sells on days 3, 4
and 6, for a profit of 9. The duplicated prices are the options to move a sale to a later, better day.}
\label{fig:regret}
\end{figure}

\begin{code}
void solve() {
    int n;
    cin >> n;
    minpq<ll> pq;  // prices at which we can still buy (or re-sell)
    ll profit = 0;
    rep(i, n) {
        ll p;
        cin >> p;
        if (!pq.empty() && pq.top() < p) {
            profit += p - pq.top();  // sell today what was cheapest
            pq.pop();
            pq.push(p);  // allows this sale to be moved to a later day
        }
        pq.push(p);  // today as an ordinary buying opportunity
    }
    print(profit);
}
\end{code}

Regret greedy is a lightweight form of the augmenting paths of min-cost flow (Chapter~\ref{ch:flows}): the ``undo'' candidates
are the reverse edges of the residual graph. When a problem looks like a flow problem but $n$ is $10^5$ or more, try to find
the greedy with undo that simulates the flow.

\begin{recognize}
Greedy is likely when the statement asks to ``maximise the number of'' or ``minimise the maximum'', the limits are large
($n = 2\cdot10^5$) with no obvious DP state, and brute-force answers on small inputs always agree with some simple rule. The
first step is almost always to sort by the right key; find the key with an adjacent exchange argument. Greedy is unlikely
when choices interact in complicated ways (then think of dynamic programming) or when the problem is a known hard problem
such as 0/1 knapsack or set cover, unless the limits reveal special structure.
\end{recognize}

\begin{summary}[Chapter summary]
\begin{itemize}
  \item After sorting, equal values are adjacent, order statistics are positions, and closest values are neighbours;
        \texttt{nth\_element} finds a median in linear time and \texttt{partial\_sort} the $k$ smallest.
  \item Counting sort and radix sort sort integers in linear time and are stable; the template sorts 64-bit integers more
        than twice as fast as \texttt{sort}.
  \item Find comparators with adjacent exchange arguments, and make sure they are strict weak orderings, ideally by
        comparing one key per element.
  \item Coordinate compression maps huge values to ranks; compress the query values too.
  \item Prove greedy algorithms by ``greedy stays ahead'' or by exchange arguments, and break wrong ones with small
        counterexamples.
  \item Heaps and ordered multisets implement greedy choices among changing candidates; regret greedy pushes undo options
        into the heap.
\end{itemize}
\end{summary}

\begin{practice}
\ppgroup{Sorting}
\begin{pplist}
  \item Distinct Numbers
  \item Apartments
  \item Ferris Wheel
  \item Concert Tickets
  \item Missing Coin Sum
  \item Collecting Numbers
\end{pplist}
\ppgroup{Greedy}
\begin{pplist}
  \item Movie Festival
  \item Stick Lengths
  \item Tasks and Deadlines
  \item Reading Books
  \item Towers
  \item Stick Divisions
\end{pplist}
\ppgroup{Greedy with data structures}
\begin{pplist}
  \item Room Allocation
  \item Movie Festival II
  \item Factory Machines
\end{pplist}
\end{practice}
